\Needspace{14\baselineskip}
\section{Responses to the comments of reviewer 1}
\label{sec:responses-reviewer-1}
\reviewercomment{The workflow is useful, but the comparison process needs a clearer explanation.}
\response{We appreciate these constructive comments. We clarified the build and comparison steps as follows.}
\begin{enumerate}
  \Needspace{5\baselineskip}
  \item \reviewercomment{Explain how the baseline is selected for each revision.}
    \response{Section IV explains that each revision names its preceding manuscript as its baseline.}
    \Needspace{8\baselineskip}
    \begin{samepage}
    \response{The comparison build performs these steps:}
    \begin{itemize}
      \item \response{Compile both manuscripts and their bibliographies.}
      \item \response{Expand their included sources and compare the complete documents.}
    \end{itemize}
    \end{samepage}
  \Needspace{5\baselineskip}
  \item \reviewercomment{Avoid manually named flattened snapshots, which can become stale.}
    \response{We removed that recommendation. Section IV now describes fixed generated filenames and automatic regeneration.}
  \Needspace{5\baselineskip}
  \item \reviewercomment{Clarify where authors find the generated PDFs and how they inspect changes.}
    \response{Section III explains that finished PDFs appear beside their manuscript sources. Section IV describes blue additions and red deletions.}
\end{enumerate}
